\chapter{What Was Missing}\label{r02:ch:gap}

This chapter walks through upstream waLBerla (\texttt{v7.2-88-g0c8ed8c9}) and upstream
\ps{} (\texttt{20d7dcf0}) on the K1 and records, layer by layer, what prevents RVV code.
The analysis is organized along the five decision points of \cref{r02:fig:build-flow}.

\section{What upstream waLBerla does on the K1}\label{r02:sec:gap-behaviour}

There are two ways to configure a CPU build, and neither yields vector code:
\begin{description}[style=nextline]
  \item[\cmakeopt{WALBERLA_OPTIMIZE_FOR_LOCALHOST=OFF} (the default)] waLBerla adds no
    \code{-march} flag. Both board compilers then default to
    \code{rv64imafdc_zicsr_zifencei}, i.e.\ \emph{without} the V extension, so not even
    the auto-vectorizer can use the vector unit. sweepgen selects
    \code{Target.GenericCPU}. Everything builds and runs, entirely scalar.
  \item[\cmakeopt{WALBERLA_OPTIMIZE_FOR_LOCALHOST=ON}] waLBerla adds
    \code{-march=native} to every compile command, and both compilers reject it:
\begin{plaincode}
gcc:   error: '-march=native': ISA string must begin with rv32 or rv64
clang: error: invalid arch name 'native', string must begin with rv32{i,e,g} or rv64{i,e,g}
\end{plaincode}
    The build fails at the first translation unit. Had it got that far, sweepgen
    would have selected \code{Target.CurrentCPU}, and upstream \ps{} would have
    resolved it to \code{GenericCPU} with the warning \emph{``Unable to determine
    available vector CPU targets for this system: unknown platform riscv64''}.
\end{description}

\section{The chain of decisions}\label{r02:sec:gap-chain}

For an x86 developer, ``build for this machine, vectorized'' works because five
independent decisions line up. On RISC-V each of them failed or was undefined:

\begin{description}[style=nextline]
  \item[\rmark{1} Compiler flags (waLBerla CMake)] On x86 the flags come from
    \code{-march=native}. On riscv64, \code{-march=native} is an error with GCC~14 and
    Clang~18. And even if it worked, it would not be enough: the K1's ISA string in
    \file{/proc/cpuinfo} (\code{rv64imafdcv_zicbom_..._zve64d_..._zvfh_...}) does
    not contain \code{zvl256b}. A compiler deriving \code{-march} from it would know
    that vectors exist, but not that they are 256 bits wide. Generated VLS code needs
    that guarantee (decision 4).
  \item[\rmark{2} Target selection (sweepgen)] sweepgen maps
    \cmakeopt{WALBERLA_OPTIMIZE_FOR_LOCALHOST} to \code{Target.CurrentCPU}, otherwise to
    \code{GenericCPU} (\cref{r02:lst:buildconfig}). This logic is
    architecture-neutral and \emph{correct} for RISC-V; it only needs \ps{} to resolve
    \code{CurrentCPU} properly. Nothing was missing here.
  \item[\rmark{3} Target resolution and code generation (\ps{})]
    \code{Target.auto_cpu()} only knows x86 (through \code{py-cpuinfo}) and ARM. There
    was no RISC-V vector target in the \code{Target} enum, no platform class, no
    intrinsic selector and no runtime header. Separately, \ps{}' own JIT compiler
    passes \code{-march=native} for \code{CurrentCPU}, which fails with GCC on RISC-V
    in the same way. (waLBerla does not use the JIT, it generates code only; but every
    \ps{} test and every Python-side experiment does.)
  \item[\rmark{4} Compiling the generated code] Generated RVV intrinsics compile only if
    the global flags enable the V extension (decision 1). VLS code with 4 lanes per
    \code{double} vector is only \emph{correct} if VLEN $\geq$ 256 (on a VLEN-128 chip,
    \code{vl = 4} exceeds VLMAX = 2 and the hardware silently processes fewer lanes).
    Some mechanism has to stop such code from being built for the wrong machine.
  \item[\rmark{5} Requesting vectorization (generator script)] sweepgen's
    \code{get_pystencils_config} never sets \code{cpu.vectorize.enable}, on any
    platform (\cref{r02:sec:buildconfig}). An application that wants explicit SIMD
    code has to say so. No waLBerla application or test in the tree does this for the
    LBM sweeps; \code{grep vectorize} over \file{sweepgen/} and \file{apps/examples/}
    finds nothing.
\end{description}

\section{Everything else that was in the way}\label{r02:sec:gap-other}

Getting from ``it generates RVV code'' to ``it builds, runs and can be measured''
uncovered more obstacles:
\begin{description}[style=nextline]
  \item[Python package versions] sweepgen's managed venv installs \ps{}, lbmpy and
    pystencils-sfg from their git \texttt{master} branches. That would replace the
    \ps{} fork, and the newest pystencils-sfg (0.2, 0.3) imports
    \code{pystencils.grids}, which does not exist in the pinned \ps{} (2026-05-21).
    pystencils-sfg 0.1.1 (2026-05-26) is the matching release. lbmpy was not installed
    at all, and the git server that sweepgen's requirements point to returned HTTP 500
    at the time; lbmpy 2.1.1 from PyPI works with the fork although it declares
    \code{pystencils~=2.1}.
  \item[Clang 18 and the v8 halo exchange] \file{GenericFieldPackInfos.hpp} declares
    \code{const ConstFieldView srcView{ field_, srcBlock };}. That deduces the template
    arguments of an \emph{alias template} (C++20 feature P1814), which GCC supports but
    Clang implements only from version 19. The board has Clang 17 and 18, so every
    application that uses the v8 halo exchange failed to compile with Clang. This is
    not RVV-specific, but it blocked the Clang half of the measurements.
  \item[No way to measure] waLBerla had no benchmark that builds the same sweep with
    and without explicit vectorization and prints a checksum. The closest examples
    (\code{DoubleShearLayer}) write VTK output and run a fixed scenario; the classic
    \code{UniformGridCPU} benchmark uses the legacy code generation path.
  \item[Board constraints] 4\,GB of RAM (the heaviest waLBerla translation units need
    about 1.3\,GB each, so at most \code{-j2}), a shared machine (everything under
    \file{/home/bananapi/saleh/}), slow builds (41 minutes for the first target of a
    fresh build directory), and a board clock about 46 hours behind the laptop.
\end{description}

\section{Not missing, but worth knowing}\label{r02:sec:gap-not}

Some parts needed \emph{no} change, which is as instructive as the gaps:
\begin{itemize}
  \item \textbf{sweepgen}: target selection is generic, and the \code{gen_config}
        argument of \cls{Sweep} and \cls{LbmBulk} already lets a script request
        vectorization. The first attempt added RVV options to sweepgen; they were
        removed again (\cref{r02:sec:design-history}).
  \item \textbf{The v8 data structures}: fields are \code{fzyx} with $x$ contiguous,
        which is what any SIMD backend needs. RVV loads only require element
        alignment, so the 16-byte line alignment on RISC-V is sufficient.
  \item \textbf{The classic SIMD module} \file{src/simd}: x86-only, but unused by
        generated code (\cref{r02:sec:simd-module}).
  \item \textbf{The legacy code generation path}: it cannot use \ps{}~2.0 at all, so
        RVV support there would mean porting the path first. Out of scope.
\end{itemize}

\Cref{r02:tab:gap-summary} summarizes the gaps and where each one was closed.

\begin{table}[htbp]
  \centering\small
  \caption{Gaps and where they were closed.}\label{r02:tab:gap-summary}
  \begin{tabularx}{\linewidth}{@{}>{\raggedright\arraybackslash}p{3.4cm}>{\raggedright\arraybackslash}X>{\raggedright\arraybackslash}p{4.1cm}@{}}
    \toprule
    \textbf{Layer} & \textbf{Gap} & \textbf{Closed by}\\
    \midrule
    \rmark{1} waLBerla CMake & \code{-march=native} rejected; no VLEN information
      & explicit \code{-march=rv64gcv_zvl256b} in \file{GNU.cmake}/\file{Clang.cmake}
      (\cref{r02:sec:wlb-cmake})\\
    \rmark{2} sweepgen & nothing & ---\\
    \rmark{3} \ps{} & no RISC-V target, platform, intrinsics, runtime header; JIT
      \code{-march=native} & \code{Target.RISCV_RVV}, \cls{RvvCpu}, \file{rvv.hpp},
      JIT flags (\cref{r02:ch:pystencils})\\
    \rmark{4} compile guard & nothing prevents VLS-256 code on a smaller VLEN &
      \texttt{\#error} in \file{rvv.hpp} unless \code{__riscv_v_min_vlen >= 256}\\
    \rmark{5} generator script & sweepgen never vectorizes & opt-in via
      \code{gen_config} in the script (\cref{r02:sec:wlb-optin})\\
    Python environment & sfg/lbmpy versions & sfg 0.1.1 + lbmpy 2.1.1 in an external
      venv, \code{PYTHONPATH} to the forks (\cref{r02:sec:env})\\
    Clang 18 & alias-template CTAD in the halo exchange & explicit type in
      \file{GenericFieldPackInfos.hpp} (\cref{r02:sec:wlb-clang})\\
    Measurement & no comparable benchmark & \file{apps/benchmarks/SweepgenLbmCPU}
      (\cref{r02:sec:app})\\
    \bottomrule
  \end{tabularx}
\end{table}
